\section{Computational scope}
\label{sec:experiments}

The repository accompanying this manuscript contains a one-dimensional
Gaussian-process example, a denoising-score-matching implementation,
a preconditioned Langevin sampler, and numerical routines for Gaussian
Hellinger distances.  These components are retained as reproducible
exploratory code.  They are not used in any proof.

In particular, a finite-batch denoising loss is not an estimator of the
path expectation $\varepsilon_N$ in \eqref{eq:drift-energy} without an
additional calibration argument.  Likewise, a log--log fit over a
small set of noise levels does not establish a posterior-contraction
exponent.  The scripts and the pilot arrays they generate should
therefore be read as implementation checks, not as empirical
verification of \Cref{thm:discrete} or \Cref{thm:contraction}.

A numerical study designed to test \Cref{thm:discrete} would need, at a
minimum, a drift architecture satisfying \Cref{ass:drift}, a direct
estimate of the Cameron--Martin energy in
\eqref{eq:drift-energy}, controlled covariance-tail truncations, and
sampling error bars.  A contraction study would additionally need an
identifiable sequence of experiments and a pre-specified loss.  Those
experiments are beyond the claims of the present paper.
